% §1245: `\deadcycles` is not an eqtb entry. TeX keeps it in a global that
% `\shipout` clears (§640) and §1025 increments as it fires `\output`:
%
%     @<Fire up the user's output routine and |return|@>=
%     begin output_active:=true;
%     incr(dead_cycles);
%     push_nest; mode:=-vmode; …
%
% so inside the output routine `\the\deadcycles` is at least 1, and it counts
% how many routines have run since the last page was shipped. `\the` reads that
% global directly.
%
% We keep the counter in a field and serve `\the` out of the name table, and
% nothing wrote one from the other — so `\the\deadcycles` answered 0 for the
% whole of every run while the counter behind it climbed. LaTeX's output
% routine reads it.
\input plain
% Nothing shipped yet, and no output routine has run: zero.
\message{<start=\the\deadcycles>}
% An output routine that ships sets it back to 0 (§640). One that does not
% leaves it standing, so a routine which simply throws the page away and
% returns is counted.
\output={\message{<in=\the\deadcycles>}\shipout\box255
  \message{<after=\the\deadcycles>}}
\vsize=20pt
\hrule height 5pt
\vfill\penalty-10000
\message{<shipped=\the\deadcycles>}
% Assigning it writes the *global*, not just the value `\the` reads (§1245's
% `set_page_int`). §1054 ends the job only when the page and the lists are
% empty **and** `dead_cycles=0`, so an assignment that reached only the
% readable copy left `\end` ejecting residual material over and over. LaTeX's
% `\enddocument` is `\deadcycles\z@\@@end`, so this is on the path out of
% every document.
\deadcycles=7
\message{<set=\the\deadcycles>}
\deadcycles=0
\message{<cleared=\the\deadcycles>}
\end
